\section{Pairwise Judge Prompt Interface}
\label{app:pairwise-prompt}
The primary pairwise evaluation used \texttt{trailtraining.llm.soft\_eval.compare\_plans}. Unlike the marker task, this interface did not expose a project-local literal prompt string. The archived reproducibility object therefore consists of: (i) the exact canonical masked JSON payload passed to the evaluator, (ii) the judge model configuration, and (iii) the returned fields \texttt{preferred}, \texttt{reasoning}, \texttt{plan\_a\_advantages}, and \texttt{plan\_b\_advantages}. The evaluator compared two source-masked plans labeled only as Plan A and Plan B. Source family, source model, plan IDs, fixture metadata, self-family indicators, and AB/BA order were not included in the judge-facing payload. The call used \texttt{SoftEvalConfig(enabled=True, reasoning\_effort=``none'', skip\_synthesis=True, parallel\_batches=False, temperature=0.0)}. Appendix file \path{appendix_artifacts/pairwise_prompt_interface.txt} records this interface, and \path{appendix_artifacts/sample_pairwise_judge_input.json} contains a complete archived pairwise payload.

\section{Marker-Level Judge Prompt}
\label{app:marker-prompt}
The marker-level evaluation used the following system prompt exactly.

\begin{lstlisting}
You are an expert running coach and careful evaluator.

You will compare two anonymized training plans, Plan A and Plan B.
The plans have been source-masked. Do not infer or mention whether either plan was generated by a model, a program, or a human.

Evaluate each marker independently. For every marker, choose:
- "plan_a" if Plan A is better on that marker,
- "plan_b" if Plan B is better on that marker,
- "tie" if they are meaningfully equivalent.

Use only the visible plan content. Do not use source labels, file names, IDs, or provenance.
Return valid JSON only. No markdown.
\end{lstlisting}

The user prompt template was:

\begin{lstlisting}
Compare Plan A and Plan B on the following markers.

MARKERS:
- plan_coherence: Overall training-plan coherence: whether the week fits together as a coherent training plan, with compatible sessions, sensible sequencing, and no internal contradictions.
- training_specificity: Specificity to the athlete and trail-running context: terrain, hills, long-run needs, race phase, and appropriate trail/endurance demands.
- load_progression: Appropriateness of training load and progression: total volume, session duration, stress distribution, and whether the plan avoids abrupt or excessive load.
- recovery_safety: Recovery, rest, and safety: rest-day placement, hard-day spacing, injury-risk control, and avoidance of unsafe combinations.
- quality_session_design: Design of quality sessions: workouts such as intervals, tempo, hills, progression runs, or long runs are purposeful, specific, and feasible.
- endurance_development: Support for aerobic/endurance development: easy volume, long-run stimulus, sustainable aerobic work, and endurance-oriented structure.
- readiness_alignment: Alignment with readiness/recovery constraints: whether the plan respects the athlete's readiness, recovery capacity, and fatigue constraints.
- clarity_actionability: Clarity and actionability: whether the athlete can understand and execute the plan without ambiguity.
- explanation_quality: Quality of explanation/rationale: whether the stated purpose and explanation are helpful, specific, and support the plan. This marker is presentation-sensitive.

Scoring guidance:
- Scores are 1 to 5, where 5 is best.
- The preferred plan should normally be the one with the higher score.
- Use tie when differences are negligible.
- Keep rationales short.

REQUIRED JSON SCHEMA:
{"markers":{"plan_coherence":{"preferred":"plan_a | plan_b | tie","plan_a_score":"integer 1-5","plan_b_score":"integer 1-5","confidence":"number 0.0-1.0","rationale":"brief reason, max 18 words"}, ...},"overall_marker_summary":"one short sentence"}

PLAN INPUT:
{compact_json(judge_input)}
\end{lstlisting}
A complete instantiated marker prompt for one archived comparison is included as \path{appendix_artifacts/sample_complete_marker_prompt.txt}.

\section{Canonical Judge-Facing Schema}
\label{app:schema}
Both sources are rendered into the same canonical judge-facing schema. The judge sees only neutral \texttt{plan\_a} and \texttt{plan\_b} objects.

\begin{lstlisting}
{
  "plan_a": {
    "judge_view": "canonical_masked_structural_v1",
    "plan": {
      "days": [
        {
          "day_index": "integer 1..7",
          "session_type": "string",
          "duration_minutes": "integer",
          "is_rest_day": "boolean",
          "is_hard_day": "boolean",
          "target_intensity": "string",
          "workout_summary": "short normalized text",
          "purpose_summary": "short normalized text"
        }
      ],
      "summary": {
        "plan_days": "integer",
        "total_minutes": "integer",
        "rest_days": "integer",
        "hard_days": "integer",
        "long_runs": "integer"
      }
    }
  },
  "plan_b": "same schema as plan_a"
}
\end{lstlisting}
The source family, judge family, source model, plan IDs, structural-score gap, self-family flag, AB/BA order, judge identity, and run index are retained in the manifest outside the prompt for analysis. The frozen artifact set is named \texttt{canonical\_masked\_scrubbed\_v1\_t000}; the judge-facing payload inside each archived JSON record uses the \texttt{canonical\_masked\_structural\_v1} view shown above.

\section{Structural-Score Feature Definitions}
\label{app:structural-features}
The structural score excludes title wording, prose richness, citations, claim attributions, data-note verbosity, rationale or explanation fields, source model, generation arm, and file names. It uses only plan structure and fixture-conditioned expectations. The scored features are:

\begin{itemize}
  \item \textbf{plan days}: actual number of plan days compared with expected length;
  \item \textbf{total minutes}: sum of daily \texttt{duration\_minutes};
  \item \textbf{active, rest, hard, quality days}: counts derived from rest and hard flags and quality session types;
  \item \textbf{long runs}: count of \texttt{session\_type == long};
  \item \textbf{maximum and mean day minutes}: daily load distribution summaries;
  \item \textbf{hard-day spacing}: minimum distance between hard days;
  \item \textbf{fixture-adjusted load, hard-day, and rest-day expectations}: ranges conditioned on athlete band, readiness, recovery capability, and race phase.
\end{itemize}

The matching distance uses weighted gaps in total minutes, number of rest days, number of hard days, number of active days, number of long runs, number of quality days, maximum day minutes, and mean day minutes. Soft calipers penalize large gaps in total minutes, maximum day minutes, rest days, hard days, long runs, and quality days. Hard calipers were not enabled in the frozen primary match because the target-cardinality design prioritized retaining 250 same-cell pairs. The machine-readable definition is included as \path{appendix_artifacts/structural_score_features.json}.

\section{Sample Judge-Facing JSON}
\label{app:sample-json}
The following excerpt shows one canonical comparison. The full JSON appears in \path{appendix_artifacts/sample_pairwise_judge_input.json}.

\begin{lstlisting}
{
  "plan_a": {
    "judge_view": "canonical_masked_structural_v1",
    "plan": {
      "days": [
        {
          "day_index": 1,
          "duration_minutes": 48,
          "is_hard_day": false,
          "is_rest_day": false,
          "purpose_summary": "Build aerobic consistency while keeping fatigue low.",
          "session_type": "easy",
          "target_intensity": "easy",
          "workout_summary": "48 min easy run at conversational effort."
        },
        {
          "day_index": 2,
          "duration_minutes": 39,
          "is_hard_day": false,
          "is_rest_day": false,
          "purpose_summary": "Build aerobic consistency while keeping fatigue low.",
          "session_type": "easy",
          "target_intensity": "easy",
          "workout_summary": "39 min easy run at conversational effort."
        }
      ],
      "summary": {
        "hard_days": 0,
        "long_runs": 0,
        "plan_days": 7,
        "rest_days": 2,
        "total_minutes": 229
      }
    }
  },
  "plan_b": {
    "judge_view": "canonical_masked_structural_v1",
    "plan": {
      "days": [
        {
          "day_index": 1,
          "duration_minutes": 38,
          "is_hard_day": false,
          "is_rest_day": false,
          "purpose_summary": "Build confidence and endurance.",
          "session_type": "easy",
          "target_intensity": "easy",
          "workout_summary": "Easy trail run focusing on form and cadence."
        },
        {
          "day_index": 2,
          "duration_minutes": 41,
          "is_hard_day": false,
          "is_rest_day": false,
          "purpose_summary": "Maintain aerobic base and hill strength.",
          "session_type": "easy",
          "target_intensity": "easy",
          "workout_summary": "Easy trail run to maintain aerobic base."
        }
      ],
      "summary": {
        "hard_days": 0,
        "long_runs": 0,
        "plan_days": 7,
        "rest_days": 1,
        "total_minutes": 321
      }
    }
  }
}
\end{lstlisting}

\section{Sample Metadata Retained Outside the Prompt}
The manifest retains analysis metadata that the judge does not see. One example is shown below; the full file appears as \path{appendix_artifacts/sample_pair_metadata_outside_prompt.json}.

\begin{lstlisting}
{
  "record_id": "pair_0001__qwen_14b_judge__r00__BA",
  "pair_id": "pair_0001",
  "fixture_id": "ab_A1__r_high__rc_high__ph_base",
  "judge": "qwen_14b_judge",
  "judge_model_family": "qwen",
  "run": 0,
  "order": "BA",
  "left_plan_role": "programmatic",
  "right_plan_role": "llm",
  "llm_plan_id": "ab_A1__r_high__rc_high__ph_base__qwen2.5_7b_instruct__src_t070__exp_t000__s004",
  "programmatic_plan_id": "ab_A1__r_high__rc_high__ph_base__prog__exp_t000__s008",
  "source_model_family": "qwen",
  "self_family_match": true,
  "structural_score_gap": 0.0,
  "structural_score_version": "structural_score_v1.0.0",
  "feature_gaps": {
    "max_day_minutes": 35.0,
    "mean_day_minutes": 13.143,
    "n_active_days": 1.0,
    "n_hard_days": 0.0,
    "n_long_runs": 0.0,
    "n_quality_days": 0.0,
    "n_rest_days": 1.0,
    "total_minutes": 92.0
  },
  "judge_view": "canonical_masked"
}
\end{lstlisting}


